\documentclass[11pt,a4paper]{article}
\usepackage{templates/modeling-notes}
\hypersetup{pdftitle={Simplex practice -- four mock labs}}
\title{\textbf{Mathematical Modeling}\\[2mm]\large Four mock labs: standard form and simplex}
\author{Personal study notes}
\date{}
\begin{document}
\maketitle

\section*{How to use these tasks}
Complete the solved demo first. Then attempt each mock lab before opening
the separate answer PDF or its reference MATLAB script.
The time budgets below are suggested practice limits. Each task is independent.
For hand work, standard form means $Sv=b$, $v\geq0$, with a maximization
objective. Convert minimization by negating the objective. For MATLAB, use
\texttt{linprog} with \texttt{dual-simplex-highs}; remember that it minimizes.
Put variable bounds into matrix rows and pass empty \texttt{lb}, \texttt{ub}.

In every submission, include variables and units, the original model, the
standard-form matrix and variable order, a starting-basis check, the requested
hand work, MATLAB arrays and output, and an interpretation in the original
variables. Check solver status, signs, constraint residuals and the objective.
Use $10^{-7}$ as the absolute comparison tolerance. Do not require identical
optimal vectors or split-variable values when the representation is nonunique.

For hand pivots use the most negative $z_j-c_j$, breaking entering ties by
smaller column index. Break leaving-ratio ties by the smaller basic-variable
index. Keep variable headings fixed. Include $c_B$, $z_j$, and $z_j-c_j$;
place slash ratios and the leaving arrow outside the RHS, as in the demo.

\section{Mock lab A: a feasible slack start (40 minutes)}
A workshop makes two divisible products. Product 1 earns 5 monetary units per
batch and uses 3 hours of assembly and 1 hour of finishing. Product 2 earns
4 per batch and uses 2 hours of each resource. There are 18 assembly hours and
10 finishing hours. Unsold inventory is not required and production may be zero.
\begin{enumerate}[leftmargin=*]
  \item Formulate the profit-maximizing LP. State the dimensions of all arrays.
  \item Add slacks $s_1,s_2$ in that order. Use the fixed column order
        $(x_1,x_2,s_1,s_2)$. Identify and verify an initial BFS.
  \item Solve by hand, showing every tableau. Explain the stopping condition.
  \item Verify both the original inequalities and standard-form equalities
        with MATLAB. Report production, unused hours and total profit.
  \item Draw the feasible polygon and mark your hand pivot path.
  \item Repeat the solver calculation with 20 assembly hours and unchanged
        finishing hours. Explain the change in production and profit.
\end{enumerate}
Suggested self-marking: model 2, standard form/start 2, pivots 3,
MATLAB checks 2, interpretation 1 (10 points total).

\clearpage
\section{Mock lab B: a start that needs Phase I (50 minutes)}
Two divisible feed components cost 2 and 3 monetary units per kg. Their masses
are $x_1,x_2\geq0$. At least 4 kg of feed is required. A processing condition
requires exactly $x_1+3x_2=6$. Minimize total cost.
\begin{enumerate}[leftmargin=*]
  \item Write the original LP. Explain why subtracting a surplus variable
        from the minimum-mass constraint does not immediately give a feasible
        slack/surplus basis.
  \item Use the original standard-form order $(x_1,x_2,s_1)$.
        Add artificials $a_1,a_2$ in that order for the hand calculation,
        giving $(x_1,x_2,s_1,a_1,a_2)$.
  \item Start with $(a_1,a_2)$ basic and maximize $-a_1-a_2$.
        Construct the indicator row from the actual basic coefficients.
        Carry out Phase I and state what its optimum proves.
  \item Remove artificial columns when justified. Recompute the original
        objective as a maximization of negative cost and determine whether
        a Phase II pivot is required.
  \item Solve the original LP with MATLAB, then solve its original
        three-variable standard form without artificial variables. Explain
        why the solver input does not require your hand Phase I variables.
  \item Independently eliminate one variable using the equality and minimize
        the resulting one-variable expression on its feasible interval.
\end{enumerate}
Suggested self-marking: model/standard form 2, Phase I 4, Phase II 2,
solver and independent check 2 (10 points total).

\section*{Preparation note}
This is an extension after the slack-start task. Read Chapter 3's initialization
section before attempting it if Phase I is unfamiliar.

\clearpage
\section{Mock lab C: recover the original variables (45 minutes)}
A continuous planning model has a signed adjustment $u$ with no sign
restriction, and a second adjustment $v$ satisfying $-1\leq v\leq2$.
Maximize the score
\[
  F(u,v)=3u+2v+5
\]
subject to
\[
  -u-v\geq-3,\qquad -u+v\leq1.
\]
The score and variables are abstract numerical quantities; no integrality is
assumed.
\begin{enumerate}[leftmargin=*]
  \item Convert the first inequality to a $\leq$ row, explaining the sign
        reversal. In the original variables, assemble all restrictions into
        $A(u,v)^T\leq b$. Do not accidentally impose $u\geq0$ or $v\geq0$.
  \item Write $u=u^+-u^-$ and $v=w-1$, with $u^+,u^-,w\geq0$.
        Derive the new RHS values, the upper bound on $w$, the objective
        coefficients, and the objective constant.
  \item Add slacks $s_1,s_2,s_3$ and use the fixed standard-form order
        $(u^+,u^-,w,s_1,s_2,s_3)$. Verify a feasible slack basis and calculate
        its meaning in the original variables.
  \item Perform the improving hand pivot(s). Recover $u,v$ and the full
        score, including the objective constant.
  \item Solve both formulations with MATLAB. Compare recovered original
        decisions and full scores. Must the solver return exactly the same
        individual values of $u^+$ and $u^-$ as your hand calculation?
  \item Give a direct upper bound on $F$ using $u+v\leq3$ and $v\geq-1$.
\end{enumerate}
Suggested self-marking: signs/domains 2, transformation 3, simplex 2,
recovery and solver comparison 2, certificate 1 (10 points total).

\clearpage
\section{Mock lab D: diagnose the outcome (40 minutes)}
For each model below, classify the mathematical outcome before interpreting
any solver output. All variables are continuous and nonnegative.
\subsection*{D1: a finite maximum}
\[
  \max(x_1+x_2),\qquad x_1+x_2\leq4,\quad x_1\leq3,\quad x_2\leq3.
\]
Convert to standard form using three slacks. Find the optimal value, give
two distinct optimal points, and describe every optimizer. Verify the model
with MATLAB. Explain why a different returned vector need not be an error.

\subsection*{D2: competing requirements}
\[
  \min(x_1+2x_2),\qquad x_1+x_2\geq5,\quad2x_1+2x_2\leq8.
\]
Give a short mathematical feasibility certificate. For an additional Phase I
check, use
\[
  x_1+x_2-s+a=5,\qquad2x_1+2x_2+t=8,
  \qquad x_1,x_2,s,t,a\geq0.
\]
Determine the minimum possible artificial sum $a$ and exhibit a point attaining
it. Compare with the solver status for the original model. Explain why no
original optimal decision vector should be reported.

\subsection*{D3: an improving direction}
\[
  \max(2x_1+x_2),\qquad x_1-x_2\leq1.
\]
Write standard form with one slack. Carry out the first pivot from the slack
basis and inspect the next entering column. Give a feasible point and a
nonnegative improving ray, verify the constraint along it, and compute the
objective as a function of its ray parameter. Compare with the solver status.

Suggested self-marking: D1 3, D2 3, D3 4 (10 points total).

\section*{Sources and editorial remarks}
All four tasks, suggested times and marking guides are invented for practice;
they are not official Lab 3 assignments. The confirmed Lab 3 topics are simplex
and standard form. Assessment of Phase I, free-variable transformations or
a particular MATLAB implementation remains unconfirmed. Tasks B--D deliberately
broaden practice. The exact lab timetable remains unconfirmed. The original
Lab 3 starter file is left untouched.
\end{document}
